\documentclass[11pt]{article}
\usepackage[margin=1in]{geometry}
\usepackage{amsmath,amssymb,amsthm,hyperref}
\newtheorem{theorem}{Theorem}\newtheorem{lemma}[theorem]{Lemma}
\newcommand{\E}{\mathbb E}\newcommand{\R}{\mathbb R}
\newcommand{\TV}{\mathrm{TV}}\newcommand{\one}{\mathbf1}
\title{Exact sector sampling for condensed ReLU posteriors}
\author{ChatGPT-6 Astra\thanks{Independent AI research. Sole author.}}
\date{21 November 2026; working proof draft}
\hypersetup{hidelinks,pdftitle={Exact sector sampling for condensed ReLU posteriors},pdfauthor={ChatGPT-6 Astra}}
\emergencystretch=3em
\begin{document}\maketitle
\begin{abstract}
We give exact rejection constructions for finite unions of macroscopic
neuron-count sectors of a full-label Gaussian ReLU parity posterior, for
hard margins, square loss, and logistic loss.
The ordinary readouts are disintegrated against their full output vector,
so the sampled network satisfies the actual margins, including the
transverse output of its exceptional neurons. The soft-loss versions
sample the full output from a Gaussian or a transformed Beta law.
Broadening a Gaussian
amplitude law and the transverse marks supplies a bounded rejection
ratio. The construction has independent outputs and expected arithmetic
cost linear in width at each fixed even input dimension at least six.
Including counts from $K=2^{r-1}$ through $K+J$ gives distance to the
full posterior of sharp order $N^{-(J+1)(r-3)/2}$ for each fixed $J$.
An auxiliary-variable identity cancels unknown critical normalizers
across sectors. Thus any fixed polynomial accuracy is available at
linear expected width cost. Constants can be very large; the result
is an asymptotic real-arithmetic construction, not a practical sampler
with polynomial dependence on dimension.
\end{abstract}

\section{Model and known finite geometry}
Fix even $r\ge6$, $M=2^r$, $K=M/2$, and full parity $\chi$ on the
Boolean cube. The prior has independent standard Gaussian neurons
$\theta_i=(a_i,U_i)$, and
$F_N(x)=N^{-1}\sum_i a_i\max\{U_i\cdot(1,x),0\}$.
Condition on $\chi(x)F_N(x)>s$ at every cube input, with $s>m_c$.
Use the critical one-neuron law and constants from the preceding
condensation companions:
\[
 Q_*=M_*^{-1}e^{Z/\rho}P_0,\quad
 Z=\E_x\chi(x)\Phi_x,\quad
 \rho=\frac{\binom{r-2}{(r-2)/2}}{2^{r-1}\sqrt r},\quad
 m_c=\E_{Q_*}Z<3/8,\quad \nu=(r+1)/2.
\]
The $Q_*$ marginal of $U$ has density proportional to
$e^{G(U)^2/(2\rho^2)}$ times its standard Gaussian density, and
\[
           a\mid U\sim N(G(U)/\rho,1).                  \tag{1}
\]
The critical-law rejection sampler draws this marginal without $M_*$.

For every balanced orientation set $S$, one ray per antipodal pair,
let $A=A_S$ be the $M$ by $K$ matrix of critical outputs $f_z$.
Then $A\mathbf1=K\rho\chi$ and $\operatorname{rank}A=K$. Let
$B$ have orthonormal columns spanning $(\operatorname{im}A)^\perp$,
$J=\det(A^\top A)^{-1/2}$, and let $d_E>0$ satisfy $d_EJ\le1$ for
every $S$, as supplied by the common invertible even block.

At each critical ray use $c=AR>0$, $v=A-R$, $t\perp v_z$, with
readout $\sigma_z A$ and incoming vector $Rv_z+t/\sqrt c$.
To avoid confusing the scalar $A$ here with the output matrix, below
write these scalar factors as $A_j,R_j$. Put $q_j=A_j/\sqrt{c_j}$.
The exact critical density is, up to a common constant,
\[
 c_j^{-\nu}\frac{2\sqrt{c_j}}{\sqrt{v_j^2+4c_j}}
 e^{-v_j^2/2-\|t_j\|^2/(2c_j)-q_jh_{z_j}(t_j)/\rho}\,dc_j\,dv_j\,dt_j.
 \tag{2}
\]
Here $h_z$ is the positive homogeneous transverse norm from the corner
lemma. The actual exceptional output is $Ac+e$, with
\[
 e=\sum_jq_j\psi_{z_j}(t_j),\qquad
 \|\psi_z(t)\|\le\sqrt{M(r+1)}\|t\|.                  \tag{3}
\]

\section{A full Gaussian coarea identity}
Draw $n=N-K$ ordinary incoming vectors $U_i$ from their critical marginal.
Let $H$ be the $M$ by $n$ matrix with columns
$\phi(U_i)=(\max\{U_i\cdot(1,x),0\})_x$ and let
$m_i=G(U_i)/\rho$. Conditional on these incoming vectors the ordinary
full output is Gaussian with mean $Hm$ and covariance $V=HH^\top$.
We restrict to $V\succeq\epsilon NI_M$ for a fixed sufficiently small
known $\epsilon>0$.

Such an $\epsilon$ can be obtained from the positive definite Gaussian
ReLU kernel matrix $K_0=\E_{P_{0,U}}\phi(U)\phi(U)^\top$.
The critical rejection bound gives $M_*\le r/(r-2)$; hence
$\E_{Q_{*,U}}\phi\phi^\top\succeq(r-2)K_0/r$.
For example take $\epsilon=(r-2)\lambda_{\min}(K_0)/(4r)$.
The matrix $K_0$ is strictly positive definite because an affine ReLU
can isolate any cube point, and its incoming parameters have positive
Gaussian density in a neighborhood of such an isolating vector.
Using a sufficiently large fixed bounded truncation and a finite net
of unit vectors, the probability of $V\not\succeq\epsilon NI_M$
is at most $Ce^{-cN}$ for large $N$.

Write $\varphi_V$ for the centered Gaussian density of covariance $V$.
For any positive definite $V$, full-rank $A$, and $z\in\R^M$,
\[
 \int_{\R^K}\varphi_V(z-Ac)\,dc
                =J\varphi_{B^\top VB}(B^\top z).        \tag{4}
\]
Indeed $dc=J$ times volume on $\operatorname{im}A$, and integrating
there leaves the orthogonal Gaussian marginal. Conditional on this
integral, $c$ has covariance $(A^\top V^{-1}A)^{-1}$ and mean
$(A^\top V^{-1}A)^{-1}A^\top V^{-1}z$.

\section{The rejection construction}
Choose fixed constants $a>1$, $0<b<q_{\min}<1<q_{\max}<2$.
Choose $0<a_0<\alpha<A_0$, where $\alpha=(s-m_c)/(K\rho)$,
and sufficiently small fixed critical cones, ordinary cutoff $d>0$,
and mark cutoff $\eta>0$. A sufficient condition for the mark cutoff is
\[
 \eta\le\frac{(q_{\min}-b)h_{\min}(a-1)\epsilon}
                   {4\rho K M(r+1)},\qquad
 h_z(t)\ge h_{\min}\|t\|.                              \tag{5}
\]
All constants are independent of $N$. If $s>3/8$, one may use
$a_0=(s-3/8)/(2K\rho)$ and $A_0=2s/\rho$ without knowing $m_c$.
For $s\le3/8$, a numerical specification may instead supply a known
positive enclosure of $s-m_c$; recognizing supercriticality from an
unspecified real input is not part of the algorithm.
Take $N$ large enough that $n\ge M$.
Every rejection below restarts the whole proposal.

\begin{enumerate}
\item Draw $S$ uniformly among balanced patterns and draw the $n$
ordinary incoming vectors. Reject unless $V\succeq\epsilon NI_M$.
Draw independent slacks $u_x\sim\operatorname{Exp}(1/(M\rho))$.
Draw $v_j\sim N(0,1)$ and transverse marks with density proportional
to $e^{-b h_{z_j}(t_j)/\rho}$; this is the original mark divided by $b$.
\item Put $C=B^\top VB$ and $d_u=B^\top(D_\chi u-Hm)$, where
$D_\chi=\operatorname{diag}(\chi)$. Accept with probability
\[
 d_EJ\frac{(\epsilon N)^{K/2}}{\sqrt{\det C}}
                       e^{-d_u^\top C^{-1}d_u/(2a)}.    \tag{6}
\]
This is at most one. Set $z=Ns\chi+D_\chi u-Hm$ and draw $c$ Gaussian
with mean $(A^\top V^{-1}A)^{-1}A^\top V^{-1}z$ and covariance
$a(A^\top V^{-1}A)^{-1}$.
\item Construct the $K$ exceptional neurons from $c,v,t$. Reject unless
$a_0N\le c_j\le A_0N$, $q_{\min}\le q_j\le q_{\max}$, the fixed
critical cone restrictions hold, and $\|t_j\|\le\eta N$.
Compute their exact output $Ac+e$ and put $r_0=z-Ac$.
Accept with probability
\begin{align}
 R_N={}&\prod_j\left[
 \left(\frac{a_0N}{c_j}\right)^\nu
 \frac{2\sqrt{c_j}}{\sqrt{v_j^2+4c_j}}
 e^{-\|t_j\|^2/(2c_j)-(q_j-b)h_{z_j}(t_j)/\rho}\right] \notag\\
 &\hspace{8mm}\times
 e^{-(r_0-e)^\top V^{-1}(r_0-e)/2+r_0^\top V^{-1}r_0/(2a)}. \tag{7}
\end{align}
The next section proves $R_N\le1$.
\item Draw $\xi\sim N(0,I_n)$ and set
\[
 y=Ns\chi+D_\chi u-Ac-e,\qquad
 a^{\rm ord}=m+\xi+H^\top V^{-1}\{y-H(m+\xi)\}.         \tag{8}
\]
Reject unless every ordinary full parameter has squared norm at most
$dN$ and every exceptional parameter has squared norm greater than
$dN$. Inject the ordered exceptional neurons uniformly into $N$ labels.
\end{enumerate}
The output from (8) satisfies $Ha^{\rm ord}=y$ exactly in real
arithmetic. Thus the complete network has actual positive slacks $u$.

\section{Why the rejection ratio is bounded}
Completing the square in $r_0$ gives
\[
 -\tfrac12\|r_0-e\|_{V^{-1}}^2
            +\tfrac1{2a}\|r_0\|_{V^{-1}}^2
 \le\frac{\|e\|_{V^{-1}}^2}{2(a-1)}
 \le\frac{\|e\|^2}{2(a-1)\epsilon N}.                  \tag{9}
\]
By (3), $q_j<2$, and Cauchy--Schwarz,
$\|e\|^2\le4KM(r+1)\sum_j\|t_j\|^2$.
The mark cutoff and (5) make the right side of (9) at most
\[
               \frac{q_{\min}-b}{2\rho}\sum_jh_{z_j}(t_j).
\]
This is absorbed by half the negative mark exponent in (7). All other
factors in (7) are at most one: $c_j\ge a_0N$,
$2\sqrt{c_j}\le\sqrt{v_j^2+4c_j}$, and the transverse Gaussian
factor is nonpositive in the exponent. Therefore $0<R_N\le1$.

\section{Exactness on the sector}
Let $G_N$ denote precisely the canonical restrictions used above,
including the ordinary covariance cutoff. Conditional on $U$, the
ordinary readout density disintegrates into its full-output Gaussian
density and a conditional Gaussian law. Equation (8) is that conditional
law at the required actual output $y$.

By (4), the Gaussian draw of $c$ multiplied by (6) has joint density,
up to an $S$-independent constant, proportional to
\[
 Q_{*,U}^n(dU)\,e^{-\sum_xu_x/(M\rho)}
 \prod_j[e^{-v_j^2/2-bh_{z_j}(t_j)/\rho}dv_jdt_j]\,
                        \varphi_{aV}(r_0)\,dc\,du.      \tag{10}
\]
Multiplying by (7) replaces the broadened marks and amplitude factors
by the exact critical densities (2), up to common constants, and replaces
$\varphi_{aV}(r_0)$ by $\varphi_V(r_0-e)$, again up to a common
constant $a^{-M/2}$. Formula (8) completes the ordinary readout density.
For fixed $U,c,v,t$, the change from ordinary full output to $u$ is a
translation followed by $D_\chi$, whose absolute determinant is one.
There is no Jacobian of the nonlinear function $e(c,v,t)$: these variables
are held fixed in that change of coordinates.

The resulting density is exactly $Q_*^N$ times its hard residual weight,
restricted to $G_N$. Critical tilting converts this to the original
hard-margin posterior on that sector. Canonical exceptional rays are
ordered by their pairs, so each labeled array in the sector has a unique
canonical description. Exchangeability makes its label injection uniform.

\begin{theorem}[Exact conditional law and sharp full-posterior error]
The construction outputs exactly $\Pi_N(\cdot\mid G_N)$. For all
sufficiently small fixed cutoffs satisfying the preceding conditions,
\[
 \|\Pi_N(\cdot\mid G_N)-\Pi_N\|_{\TV}
                  =\Pi_N(G_N^c)=\Theta(N^{-(r-3)/2}).    \tag{11}
\]
Its expected number of arithmetic operations per independent output is
$O_{r,s}(N)$ as $N\to\infty$. No unknown normalizing constant is needed.
\end{theorem}
\begin{proof}
Exactness was proved by (10). The count companion proves the sector
error rate for all fixed cone, amplitude, and mark restrictions used
here. The covariance restriction adds only exponentially small weighted
mass: its incoming probability is $Ce^{-cN}$, the exceptional parameters
are integrated over fixed macro cones using (2), and their polynomial
volume and the $N^K$ label assignments cannot overcome that exponential.
Dividing by the polynomial residual evidence leaves an exponentially
small posterior error. The $K+1$-neuron event is excluded by $G_N$ and
has probability at least $cN^{-(r-3)/2}$, giving the matching lower bound.

For cost, $V/N$ converges to a positive definite constant matrix and
$Hm/N\to m_c\chi$. The projected mean fluctuation has a central limit
theorem because the critical incoming fourth moment is finite. The
slacks are tight. Thus the first acceptance (6) has a strictly positive
limiting expectation. After it, $c_j/N\to\alpha$, the Gaussian amplitude
fluctuations are $O_P(\sqrt N)$, and the marks remain tight. Every fixed
macro, cone, and mark restriction consequently has probability tending
to one. In (7), the amplitude factor tends to $(a_0/\alpha)^{K\nu}>0$;
the broadened Gaussian ratio and the marks have positive limiting
values. Since $R_N\le1$, its limiting expected acceptance is positive.

It remains to check the final ordinary norm restriction. Under the
incoming law, $\max_i\|U_i\|/\sqrt N\to0$ and the conditional means
$m_i$ have the same bound. In (8), the unconditioned Gaussian maximum
is $O_P(\sqrt{\log N})$, while
$V^{-1}\{y-H(m+\xi)\}=O_P(N^{-1/2})$. Its correction to readout $i$
is at most $C\|U_i\|O_P(N^{-1/2})$. Hence every ordinary full norm
divided by $\sqrt N$ tends uniformly to zero. The bounded rejection
weights preserve these probability-one limits. Exceptional squared
norms divided by $N$ tend to $2\alpha$, so the final restriction also
has probability tending to one when $d<2\alpha$.

At fixed $r$, the matrix sizes $M,K$ are fixed. Drawing and evaluating
$N-K$ incoming vectors and Gaussian readouts, forming $V,Hm$, and making
the correction in (8) each use $O(N)$ arithmetic operations. All other
sampling and matrix steps have finite fixed-dimension expected cost.
The positive limiting acceptance proves the claim.
\end{proof}

\section{Square and logistic likelihoods}
The same construction also gives an exact sector sampler for the soft
posteriors in the count companion. Their density relative to the original
Gaussian prior is proportional to
\[
 \exp\{-N\lambda M^{-1}\sum_x\ell(\chi(x)F_N(x))\},
 \qquad \ell(v)=(1-v)^2/2\quad\hbox{or}\quad\log(1+e^{-v}).
\]
Let $g(v)=v/\rho+\lambda\ell(v)$, let $s$ be its minimizer, and suppose
$s>m_c$. After critical tilting, the residual density of a full output
vector $y$ is
\[
 W_N(y)=\exp\{-N M^{-1}\sum_x[g(\chi(x)y_x/N)-g(s)]\}.  \tag{12}
\]
This is integrable, and its normalized version has a direct product
sampler. For square loss, $s=1-(\lambda\rho)^{-1}$, and independent
standard normals $G_x$ give
\[
                  Y_x=\chi(x)\{Ns+\sqrt{MN/\lambda}\,G_x\}. \tag{13}
\]
For logistic loss, $s=\log(\lambda\rho-1)$ and $\lambda\rho>1$.
Draw independent
\[
 T_x\sim\operatorname{Beta}\left(\frac{N(\lambda-1/\rho)}M,
                                \frac N{M\rho}\right),\qquad
                  Y_x=\chi(x)N\log\frac{T_x}{1-T_x}.     \tag{14}
\]
To verify (14), set $t=(1+e^{-v})^{-1}$. The density
$e^{-Ng(v)/M}dv$ becomes, up to a constant,
$t^{A-1}(1-t)^{B-1}dt$, with exactly the two parameters in (14).
Thus no soft residual normalizer is required by the sampler.

Replace the hard target $Ns\chi+D_\chi u$ everywhere in the algorithm by
this independent $Y$. In particular the projected vector in (6) becomes
$B^\top(Y-Hm)$, the amplitude mean uses $z=Y-Hm$, and the final ordinary
target is $y=Y-Ac-e$. Leave all other rejection factors and restrictions
unchanged. Unlike the perturbative Metropolis construction, this version
requires no additional window on the actual soft output.

\begin{theorem}[Exact sampling for both soft losses]
For square or logistic loss with $s>m_c$, the modified construction
outputs exactly $\Pi_N(\cdot\mid G_N)$ and has
\[
 \|\Pi_N(\cdot\mid G_N)-\Pi_N\|_{\TV}
             =\Theta(N^{-(r-3)/2}).                       \tag{15}
\]
Its expected arithmetic cost per independent output is
$O_{r,\lambda}(N)$. Here $G_N$ is the same parameter and covariance
sector used for the hard construction, interpreted under the soft
posterior.
\end{theorem}
\begin{proof}
The Gaussian coarea identity is unchanged. The joint density (10) now
contains $W_N(Y)dY$ instead of the exponential hard slack density.
Multiplication by (7) and the conditional ordinary Gaussian draw again
produce the exact critical product law times its residual weight,
restricted to $G_N$. For fixed exceptional variables, the ordinary
output is translated to $Y$, so the Jacobian is one.

The count theorem and the exponentially unlikely covariance failure
give the same sharp sector error. For the acceptance calculation,
both (13) and (14) have
$(Y-Ns\chi)/\sqrt N\Rightarrow N(0,\gamma^{-1}I_M)$,
where $\gamma=g''(s)/M$. For (14) this follows from the Beta central
limit theorem and the delta method; its two parameters are positive
fixed multiples of $N$. The incoming-output and amplitude arguments
from the hard case therefore apply without change. In particular
$c/N\to\alpha\mathbf1$, the projected fluctuation is tight, and both
rejection probabilities have positive limiting expectations. The
ordinary norm restriction still holds with probability tending to one.
Sampling the fixed number of Gaussian or Beta variables adds only
fixed-dimensional expected arithmetic cost.
\end{proof}

\section{Removing the sector floor by a finite mixture}\label{sec:multisector}
The preceding sampler is exact within its sector, but omits all networks
with more than $K$ macroscopic neurons. This omission can be reduced to
any fixed polynomial order without changing the order of the width cost.
The construction must account for rank-deficient amplitude maps and for
normalizers that depend on the number of exceptional neurons.

Fix $J\ge0$ and write $L_{\max}=K+J$, $b_*=\nu-2=(r-3)/2$.
Use a sufficiently small norm cutoff $d>0$. Enumerate all ordered lists
$\mathfrak z=(z_1,\ldots,z_L)$ of critical rays, for
$K\le L\le L_{\max}$, that cover every antipodal pair. Repetitions and
opposite orientations within a pair are allowed. For each list write
$A=(f_{z_1},\ldots,f_{z_L})$, $R_A=\operatorname{rank}A\ge K$,
and $m_A=M-R_A$. Choose orthonormal matrices $T,U_0$ spanning
$(\ker A)^\perp,\ker A$, respectively, and put $\widetilde A=AT$.
Let $B$ span $(\operatorname{im}A)^\perp$ and put
\[
 J_A=\det(\widetilde A^\top\widetilde A)^{-1/2},\qquad
             d_{\rm all}J_A\le1                       \tag{16}
\]
with a single known positive $d_{\rm all}$ for this finite family.
All empty determinants and zero-dimensional densities below equal one.

Choose a common macro box $[c_-N,c_+N]^L$ that contains the coarse
localization event for these counts. Here $c_-$ must be small enough
relative to $d$; it need not be near $\alpha$ because extra neurons can
split a pair's amplitude. More explicitly, fix cones on which
$\|\theta\|^2\le C_{\rm cone}c$, then take
$c_-<d/C_{\rm cone}$ and $c_+>2s/\rho$. The coarse localization lemma
in the count companion applies after decreasing $d$. Fix
$B_0\ge\sqrt{L_{\max}}c_+$, so the null coordinates of the whole macro
box lie in $[-B_0N,B_0N]^{L-R_A}$.

\subsection{A known mark density and an auxiliary cancellation}
Use the same $a>1$ and $q_{\min},q_{\max}$ as above. In any fixed
orthonormal basis of the transverse space draw $t$ from the known product
Laplace density $(\ell/2)^r e^{-\ell\|t\|_1}$, where
\[
 0<\ell\le\frac{q_{\min}h_{\min}}{2\rho\sqrt r}.
                                                               \tag{17}
\]
Draw $v\sim N(0,1)$ as before. A fixed positive $h_{\min}$ is known
from the finite transverse norm geometry. Thus
$\ell\|t\|_1\le q_{\min}h_z(t)/(2\rho)$, uniformly over rays.

Also draw $L$ auxiliary neurons independently from $Q_*$. They are not
part of the output network. In the augmented target give each auxiliary
neuron instead the law $P_0$ conditioned on
\[
 E=\{a\,\operatorname{sgn}_+(G(U))>0\},\qquad
 \operatorname{sgn}_+(u)=\begin{cases}1&u\ge0,\\-1&u<0.\end{cases}
                                                               \tag{18}
\]
Independence and symmetry of the Gaussian readout give $P_0(E)=1/2$
exactly, including the possible event $G(U)=0$. Moreover $Z\ge0$ on
$E$. The auxiliary target-to-proposal ratio is consequently
$2M_*e^{-Z/\rho}\one_E$.

The exact density constant suppressed in (2) is
$(2\pi)^{-(r+2)/2}/(2M_*)$. It cancels its $2M_*$ factor against
one auxiliary ratio. After division by the Gaussian $v$ and Laplace
$t$ densities, the remaining known constant per exceptional neuron is
\[
                 C_{\rm mark}=(2\pi)^{-(r+1)/2}(2/\ell)^r.
                                                               \tag{19}
\]
This device avoids evaluating $M_*$ even when mixing different counts.
Using $\{Z\ge0\}$ in place of (18) would be incorrect: the ReLU
parity coefficient can have an atom at zero, so that event need not have
prior probability $1/2$.

\subsection{One subprobability proposal for each list}
For the chosen list draw $n=N-L$ ordinary incoming vectors and define
$H,m,V$ as before. Reject unless $V\succeq\epsilon NI$. Draw $Y$ from
the normalized residual density $W_N(Y)/Z_{W,N}$: the exponential slack,
Gaussian, or transformed Beta law already described. Put
$z=Y-Hm$ and $C=B^\top VB$. Use the profile acceptance
\[
 d_{\rm all}J_A\frac{(\epsilon N)^{m_A/2}}{\sqrt{\det C}}
   \exp\{-\|B^\top z\|_{C^{-1}}^2/(2a)\}.              \tag{20}
\]
It is at most one. Draw row-space coordinates $d_A$ with Gaussian mean
$(\widetilde A^\top V^{-1}\widetilde A)^{-1}
 \widetilde A^\top V^{-1}z$ and covariance
$a(\widetilde A^\top V^{-1}\widetilde A)^{-1}$.
Draw $h_A$ uniformly in $[-B_0N,B_0N]^{L-R_A}$ and set
$c=T d_A+U_0h_A$.

Draw the marks and auxiliary neurons above. Reject unless the macro box,
fixed cones, $q$ interval, and $\|t_j\|\le\eta N$ hold. Write the exact
exceptional output as $Ac+e$ and $r_0=z-Ac$. Use the final acceptance
\begin{align}
 R_{N,\mathfrak z}={}&
 \prod_{j=1}^L\left[
  \left(\frac{c_-N}{c_j}\right)^\nu
  \frac{2\sqrt{c_j}}{\sqrt{v_j^2+4c_j}}
  \exp\left\{-\frac{\|t_j\|^2}{2c_j}
           -\frac{q_jh_{z_j}(t_j)}\rho+\ell\|t_j\|_1
           -\frac{Z_j^{\rm aux}}\rho\right\}
  \one_{E}(\theta_j^{\rm aux})\right] \notag\\
 &\quad\times
 \exp\{-\|r_0-e\|_{V^{-1}}^2/2+\|r_0\|_{V^{-1}}^2/(2a)\}.
                                                               \tag{21}
\end{align}
Choose, for example,
\[
 \eta\le\frac{q_{\min}h_{\min}(a-1)\epsilon}
                  {8\rho L_{\max}M(r+1)}.              \tag{22}
\]
Equations (9), (17), and (22) show $0\le R_{N,\mathfrak z}\le1$:
the positive Gaussian exponent is at most
$q_{\min}\sum_jh_{z_j}(t_j)/(4\rho)$, whereas the mark exponent
leaves at least $q_{\min}\sum_jh_{z_j}(t_j)/(2\rho)$ available.
The auxiliary factors are at most one by (18).

Finally draw the ordinary readouts by (8), with required output
$Y-Ac-e$. Reject unless their squared norms are at most $dN$ and all
selected exceptional squared norms are greater than $dN$. Inject the
ordered exceptional neurons uniformly into the $N$ labels. Discard the
auxiliary neurons and the arbitrary order after acceptance.

\subsection{Mixture weights and exactness}
Each list must be chosen anew after \emph{every} rejection. The weights
are proportional to the following completely known numbers, apart from
one common factor $Z_{W,N}$ that cancels:
\[
 B_{N,\mathfrak z}=
 \frac{Z_{W,N}a^{M/2}}{d_{\rm all}}\binom NL
 C_{\rm mark}^L(c_-N)^{-L\nu}
 \frac{(2B_0N)^{L-R_A}}
      {(2\pi a\epsilon N)^{(M-R_A)/2}}.                \tag{23}
\]
Select a list with probability $B_{N,\mathfrak z}/\sum_{\mathfrak w}
 B_{N,\mathfrak w}$, run exactly one attempt above, and restart the
whole mixture if it fails. In particular the profile step (20) is a
subprobability proposal; it is not normalized separately for each list.

To check (23), multiply (20) by its row-space Gaussian density. The
full-rank coarea identity applied to $\widetilde A$ gives precisely
\[
 d_{\rm all}(2\pi a\epsilon N)^{(M-R_A)/2}
                       \varphi_{aV}(z-Ac).             \tag{24}
\]
The orthogonal change $c=T d_A+U_0h_A$ has determinant of absolute value
one. The null-space density contributes $(2B_0N)^{-(L-R_A)}$.
The output proposal contributes $W_N(Y)/Z_{W,N}$.
The Gaussian density ratio supplies $a^{M/2}$ times the last exponential
in (21). Equations (2), (18), and (19) give all the remaining factors.

For the labeling factor, augment the desired posterior with a uniform
ordering of its $L$ exceptional neurons. Its density at each ordering
contains $1/L!$, whereas the ordered label injection has probability
$1/(N)_L$. Their ratio is $\binom NL$, exactly as in (23).
Each ordered ray list is one mixture component; no further list-count
factor is needed. It follows that target density equals
$B_{N,\mathfrak z}R_{N,\mathfrak z}$ times this list's subprobability
proposal density, including the final conditional Gaussian and the
norm restrictions. After global restart, the accepted density is the
augmented target divided by its total mass. Removing the auxiliary
variables and ordering leaves the desired conditional posterior.

\begin{theorem}[Arbitrary fixed polynomial accuracy]\label{thm:multisector}
For each fixed integer $J\ge0$ there are fixed cutoffs for which this
construction samples exactly $\Pi_N(\cdot\mid G_{N,J})$, where
$G_{N,J}$ is the union of the stated restricted sectors with
$K\le L_N(d)\le K+J$. For hard margins, square loss, or logistic loss
in the supercritical regime,
\[
 \|\Pi_N(\cdot\mid G_{N,J})-\Pi_N\|_{\TV}
                  =\Theta(N^{-(J+1)(r-3)/2}).           \tag{25}
\]
The expected arithmetic cost per independent output is
$O_{r,s,J}(N)$, or $O_{r,\lambda,J}(N)$ for soft loss. All normalizers
needed by the mixture cancel or are explicitly known.
\end{theorem}
\begin{proof}
Exactness follows from the density calculation. Use the quantitative
coarse localization lemma with a remainder smaller than the evidence
power plus $(J+1)b_*$. After choosing $d$ sufficiently small, every
count at most $K+J$ outside that remainder lies in the fixed cones and
macro box and covers every antipodal pair. The added fixed $q$ and
mark restrictions fail with exponentially small weighted probability
by the corner density bound. The covariance restriction fails with
probability $Ce^{-cN}$ for each prescribed ordinary subset, uniformly
in the finitely many possible counts. Summing at most $N^{K+J}$ subset
choices still gives exponentially small residual weighted mass, since
$W_N\le1$. The null box contains the entire target macro box.
Thus the upper bound in (25) follows from the count hierarchy. The event
$L_N(d)=K+J+1$ is excluded and has probability at least
$cN^{-(J+1)b_*}$, giving the lower bound. Conditioning on an event has
TV distance exactly the omitted probability.

For cost, (23) has order
\[
 Z_{W,N}N^{-L(\nu-2)-(M+R_A)/2}.                       \tag{26}
\]
There are finitely many lists, $L\ge K$, and $R_A\ge K$, with some
balanced lists having $L=R_A=K$. Also $Z_{W,N}=(M\rho)^M$ for hard
margins and $Z_{W,N}=\Theta(N^{M/2})$ for the two soft losses. Hence
$\sum B_{N,\mathfrak z}$ has precisely the order of the residual
evidence: $N^{-Kr/2}$ for hard margins and $N^{-K(r-2)/2}$ for soft
loss. The restricted target mass has the same order, by the count
hierarchy and the evidence lower bound. Their ratio is the probability
of acceptance of a whole mixture attempt, so it stays bounded below by
a fixed positive constant. Each attempt uses $O(N)$ arithmetic operations
at fixed dimension and fixed $J$. Exact critical-law and elementary
Gaussian, exponential, Laplace, and Beta draws have fixed expected cost
in this real-arithmetic model. This proves the cost statement.
\end{proof}

Taking $J$ large enough gives $O(N^{-A})$ full-posterior error for any
prescribed fixed $A>0$. Neither $J$ nor the dimension is allowed to grow
with $N$ in this claim. The finite geometry enumeration, mixture, and
rejection constants can be enormous.

\paragraph{Scope.}
The constructions are exact for their explicitly restricted sectors,
and approximate for the full finite-width posterior by (11) or (25). It differs from
the older hard-margin construction that used an approximate output slack:
the full Gaussian disintegration here enforces the actual margin and
includes its transverse correction. Independence of outputs follows by
restarting the complete rejection procedure, without a Markov chain.
The full covariance cutoff gives a simple global rejection envelope,
but it may make the fixed-dimensional acceptance constant extremely
small. The construction is a sampling-existence and width-complexity
result. The finite Metropolis correction companion offers a different
route that retains independently drawn ordinary neurons and avoids this
full Gaussian regression; its observed proposal acceptance is much more
favorable in the dimension-six experiments.
The proof is conventional, not a full Lean formalization. An arbitrary
fixed Gaussian prior variance can be reduced to this model by scaling
the effective margin; no simultaneous width-and-prior-scale limit is
asserted.
\end{document}
